\section{Data and walk-forward design}
\label{sec:design}

\begin{figure}[t]
\centering
\begin{tikzpicture}[
  node distance=4mm and 6mm,
  box/.style={draw, rounded corners=2pt, align=center, font=\footnotesize,
              minimum height=7mm, inner sep=3pt, fill=gray!6},
  arr/.style={-{Latex[length=2mm]}, semithick}]
  \node[box] (ret) {S\&P 500 daily returns\\(data through $t-1$)};
  \node[box, below left=5mm and -6mm of ret] (corr) {Rolling correlation\\$W=252$ days};
  \node[box, below=of corr] (eig) {Eigen-decomposition\\+ MP edge $\lambda_+$};
  \node[box, below=of eig] (modes) {Candidate modes\\(Pool A / Pool B)};
  \node[box, below=of modes] (mf) {Mode features\\(Table~\ref{tab:features})};
  \node[box, below right=5mm and -6mm of ret] (mk) {Market features\\vol, corr, dispersion};
  \node[box, below=of mk] (hmm) {2-state HMM,\\filtered $P(\text{stress})$};
  \node[box, below=10mm of mf, xshift=22mm] (ml) {Walk-forward ML\\$\hat P(\text{mode reverts})$};
  \node[box, below=of ml] (port) {Top third; risk-parity mode portfolio;\\regime exposure overlay};
  \node[box, below=of port] (wf) {Hold $H=21$ days, then re-estimate};
  \draw[arr] (ret) -- (corr);
  \draw[arr] (ret) -- (mk);
  \draw[arr] (corr) -- (eig);
  \draw[arr] (eig) -- (modes);
  \draw[arr] (modes) -- (mf);
  \draw[arr] (mk) -- (hmm);
  \draw[arr] (mf) -- (ml);
  \draw[arr] (hmm) -- (ml);
  \draw[arr] (hmm.east) -- ++(4mm,0) |- (port.east);
  \draw[arr] (ml) -- (port);
  \draw[arr] (port) -- (wf);
  \draw[arr, dashed] (wf.west) -- ++(-30mm,0) |- (ret.west);
\end{tikzpicture}
\caption{Pipeline. Each arrow uses only information available at the
close before the episode starts. The dashed arrow marks the monthly
walk-forward step.}
\label{fig:pipeline}
\end{figure}


\paragraph{Data.} We use daily dividend- and split-adjusted closing
prices from Yahoo Finance for the \DataAssets{} current S\&P~500
constituents with a complete price history from \DataStart{} to
\DataEnd{} (\DataDays{} trading days of log returns). Restricting to
current constituents with full histories introduces survivorship bias
(Section~\ref{sec:limitations}). Interior gaps of up to five days are
forward-filled; no other cleaning is applied.

\paragraph{Episodes.} \Cref{fig:pipeline} summarizes the pipeline. The walk-forward is organized in \emph{episodes}.
At re-estimation day $t$, the correlation matrix is estimated on the
trailing window $[t-W,t)$ with $W=252$. Its eigenvectors and the asset
volatilities $\sigma_i$ are frozen and traded over the holding period
$[t,t+H)$ with $H=21$. The holding period equals the re-estimation step,
so the episodes tile the sample without overlap: there are
\NEpisodes{} episodes from January~2011. Everything recorded about the
estimation window (features) uses data strictly before $t$, and
everything recorded about the holding period (labels, PnL) uses data
from $t$ on. Because episode $i$'s label is fully realized when episode
$i+1$ begins, a model trained on episodes $<j$ can score episode $j$
with no overlap between training labels and test data.

\paragraph{Signal.} On each trading day $d$ of an episode, mode $k$'s
position is the negative, clipped $z$-score of its cumulative price
relative to the previous $L=60$ days,
\begin{equation}
  s_{k,d} = -\operatorname{clip}\!\left(
    \frac{X_{k,d-1}-\bar X_{k}^{(L)}}{\operatorname{sd}(X_{k})^{(L)}},\,-3,\,3\right),
  \label{eq:signal}
\end{equation}
where the mean and standard deviation are taken over days $d-L,\dots,d-1$.
This bets that a displaced normal coordinate is pulled back toward its
recent equilibrium. The original specification in the code repository,
a $z$-score of the daily mode \emph{return} over 20 days, is a
short-term reversal bet; we report it for comparison.

\paragraph{Label.} The standalone forward PnL of mode $k$ in an episode
is $\pi_k=\sum_{d=t}^{t+H-1}s_{k,d}\,r_{k,d}$. The binary label is
$y_k=\mathbf 1\{\pi_k>0\}$: did trading the mode's reversion signal make
money over the next month? This ties ``mean reversion'' directly to the
quantity a strategy can monetize. As a purely statistical check we
also report the out-of-sample lag-1 autocorrelation of mode returns.

\paragraph{Candidate pools.} Mode~0 is the market mode in every window
(its overlap with the normalized equal-weight vector is
\FullOverlap{} in the full sample) and is never a candidate. We fixed
\emph{Pool~A}, modes $1$--$30$, before seeing any out-of-sample
result. It contains all signal modes ($n_s$ ranges from 3 to 11 across
episodes) plus the first part of the noise bulk, and so lets the models
learn whether the MP edge matters. \emph{Pool~B}, every non-degenerate
mode ($1$--$249$), was added after the spectrum-wide test in
Section~\ref{sec:spectrum}; results on Pool~B should be read with that
in mind.

\section{Predicting which modes revert}
\label{sec:ml}

\subsection{Features}

Every feature is computed from the estimation window only.
Table~\ref{tab:features} lists them by group. The groups follow the
mechanisms discussed above: eigenvalue position (where the mode sits
relative to the MP edge and its neighbours), estimated speed of
reversion, serial dependence, eigenvector stability (whether the mode
is the ``same'' portfolio from one month to the next), volatility
regime of the mode, and portfolio concentration (whether the mode is a
broad combination or a near-pair). We add the mode's trailing
in-window reversion performance and, for the regime-aware models, the
market-level state described in Section~\ref{sec:regimes}.

\begin{table}[t]
\centering\small
\caption{Mode-level and market-level predictors. All are computed on the
estimation window $[t-W,t)$ or, for regime variables, from data through
the close of day $t-1$.}
\label{tab:features}
\begin{tabular}{@{}llp{8.4cm}@{}}
\toprule
Group & Feature & Definition \\
\midrule
Eigenvalue & \texttt{mode\_index} & rank $k$ \\
separation & \texttt{lam\_over\_mp} & $\lambda_k/\lambda_+$ \\
 & \texttt{log\_lambda} & $\log\lambda_k$ \\
 & \texttt{eig\_gap} & $\min(\lambda_{k-1}-\lambda_k,\ \lambda_k-\lambda_{k+1})/\lambda_k$ \\
Half-life & \texttt{log\_half\_life} & $\log$ of $\ln2/(-\ln b)$, $b$ the AR(1) slope of $X_k$ over the window (capped at 252 days) \\
Autocorrelation & \texttt{ac1} & lag-1 autocorrelation of $r_k$ \\
 & \texttt{vr5} & Lo--MacKinlay variance ratio VR(5) of $r_k$ \citep{lomackinlay1988} \\
Stability & \texttt{stability} & $\max_j |v_k^{(t)\top}v_j^{(t-H)}|$ over the previous episode's modes \\
Volatility & \texttt{vol\_ratio} & 21-day volatility of $r_k$ divided by $\sqrt{\lambda_k}$ \\
Concentration & \texttt{participation} & $1/(N\sum_i v_{ik}^4)$, the effective fraction of assets in the mode \\
 & \texttt{max\_weight} & largest $|v_{ik}/\sigma_i|$ as a share of $\sum_i|v_{ik}/\sigma_i|$ \\
Trailing & \texttt{trail\_rev\_sharpe} & Sharpe ratio of the mode's signal \eqref{eq:signal} over the last 63 window days \\
\midrule
Regime & \texttt{p\_stress} & filtered HMM stress probability \\
 & \texttt{log\_vol} & log 21-day realized volatility of the equal-weight market \\
 & \texttt{avg\_corr} & 63-day average pairwise correlation \\
 & \texttt{absorption} & $\lambda_0/N$, share of variance in the market mode \citep{kritzman2011} \\
\bottomrule
\end{tabular}
\end{table}

\subsection{Models and protocol}

We compare an $\ell_2$-regularized logistic regression ($C=0.1$, features
standardized), histogram gradient boosting (depth~3, 150 iterations,
learning rate 0.05, minimum leaf size 50), and a random forest (300
trees, minimum leaf size 25). Hyperparameters were set once to
conservative values and not tuned. We pre-specified the logistic model
as primary because the signal-to-noise ratio of the problem is low.
Models are trained on an expanding window of all (episode, mode) rows
from episodes $<j$ and score episode~$j$. The first prediction is made
after \MinTrain{} episodes (about two years), so the out-of-sample
period starts on \OOSStart{}. The logistic model is re-fit every
episode and the tree ensembles every third episode. We report the
out-of-sample AUC with 95\% confidence intervals from a bootstrap that
resamples whole episodes, which respects the cross-sectional
dependence among modes in the same month.
